\documentclass[12pt,a4paper,reqno]{amsart}

% language
\usepackage[greek,english]{babel}
\usepackage[utf8]{inputenc}
\usepackage{alphabeta}

% change default names to greek
\addto\captionsenglish{
    \renewcommand{\contentsname}{Περιεχόμενα}
    \renewcommand{\refname}{Βιβλιογραφία}
    \renewcommand{\datename}{Ημερομηνία:}
    \renewcommand{\urladdrname}{Ιστοσελίδα}
}

% math 
\usepackage{amsmath,amsthm,amssymb,amscd}

% font
\usepackage[cal=euler]{mathalfa}
\usepackage{libertinus-type1}
% \usepackage{txfonts} % for upright greek letters
\usepackage{bm} % for bold symbols
\usepackage{bbm} % for the simply-looking bb symbols

% miscellaneous 
\usepackage{changepage} %for indenting environments
\usepackage{csquotes} % example: \textcquote{}
\usepackage{blkarray}
\setcounter{MaxMatrixCols}{20} % default for pmatrix is 10!!
\usepackage{ytableau}
\usepackage{array} %needed to increase the vertical length in a tabular

% drawing
\usepackage{tikz,tikz-cd}
\usetikzlibrary{shapes.misc, patterns, matrix, calc, intersections,positioning,arrows.meta}
\usepackage{graphics,graphicx}
\usepackage{float} % provides enhanced control and customization options for floating objects such as figures and tables

% colors
\usepackage{xcolor}
\definecolor{darkcandyapplered}{rgb}{0.64, 0.0, 0.0}
\definecolor{midnightblue}{rgb}{0.1, 0.1, 0.44}
\definecolor{mylightblue}{HTML}{336699}
\definecolor{burntorange}{rgb}{0.8, 0.33, 0.0}
\definecolor{iceberg}{rgb}{0.44, 0.65, 0.82}
\definecolor{applegreen}{rgb}{0.55, 0.71, 0.0}
\definecolor{canaryyellow}{rgb}{1.0, 0.94, 0.0}
\definecolor{lightgray}{rgb}{0.83, 0.83, 0.83}

% hrefs
\usepackage{hyperref}
\usepackage[noabbrev,capitalize]{cleveref}
\hypersetup{
    pdftoolbar=true,        
    pdfmenubar=true,        
    pdffitwindow=false,     
    pdfstartview={FitH},  % fits the width of the page to the window
    pdftitle={},
    pdfauthor={},
    pdfsubject={},
    pdfkeywords={},
    pdfnewwindow=true,  % links in new window
    colorlinks=true,  % false: boxed links; true: colored links
    linkcolor=darkcandyapplered,   % color of internal links
    citecolor=midnightblue,  % color of links to bibliography
    urlcolor=cyan,  % color of external links
    linktocpage=true  % changes the links from the section body to the page number
    }

% geometry
\textwidth=16cm 
\textheight=21cm 
\hoffset=-55pt 
\footskip=25pt

% thm envs (you might need to change the path)
\theoremstyle{definition}
\newtheorem{exercise}{Άσκηση}
\newtheorem{solution}{Λύση}
\newtheorem*{remark}{Παρατήρηση}
\newtheorem*{example}{Παράδειγμα}

% math ops 
% fields
\newcommand{\NN}{\mathbbmss{N}} 
\newcommand{\ZZ}{\mathbbmss{Z}} 
\newcommand{\QQ}{\mathbbmss{Q}} 
\newcommand{\RR}{\mathbbmss{R}} 
\newcommand{\CC}{\mathbbmss{C}} 
\newcommand{\KK}{\mathbbmss{K}} 
\newcommand{\FF}{\mathbbmss{F}} 
\newcommand{\HH}{\mathbbmss{H}} 

% symmetric group
\newcommand{\fS}{\mathfrak{S}}  
\newcommand{\fp}{\mathfrak{p}}  
\newcommand{\fm}{\mathfrak{m}}  

% calligraphic 
\newcommand{\aA}{\mathcal{A}} 
\newcommand{\bB}{\mathcal{B}}
\newcommand{\cC}{\mathcal{C}}
\newcommand{\dD}{\mathcal{D}}
\newcommand{\eE}{\mathcal{E}}
\newcommand{\fF}{\mathcal{F}}
\newcommand{\hH}{\mathcal{H}}
\newcommand{\iI}{\mathcal{I}}
\newcommand{\lL}{\mathcal{L}}
\newcommand{\nN}{\mathcal{N}}
\newcommand{\oO}{\mathcal{O}}
\newcommand{\pP}{\mathcal{P}}
\newcommand{\sS}{\mathcal{S}}
\newcommand{\mM}{\mathcal{M}}
\newcommand{\uU}{\mathcal{U}}

% bold
\newcommand{\bfa}{\mathbf{a}}
\newcommand{\bfe}{\mathbf{e}}
\newcommand{\bfF}{\pmb{F}}
\newcommand{\bfR}{\pmb{R}}
\newcommand{\bfv}{\mathbf{v}}
%\newcommand{\bfx}{\bm{x}}
%\newcommand{\bfx}{\mathbf{x}} 
\newcommand{\bfx}{\pmb{x}}
\newcommand{\bfX}{\pmb{X}}
\newcommand{\bfy}{\pmb{y}}
\newcommand{\bfz}{\pmb{z}}

% roman
\newcommand{\rmA}{\mathrm{A}}
\newcommand{\rmB}{\mathrm{B}}
\newcommand{\rmC}{\mathrm{C}}
\newcommand{\rmD}{\mathrm{D}} 
\newcommand{\rmH}{\mathrm{H}} 
\newcommand{\rmh}{\mathrm{h}} 
\newcommand{\rmI}{\mathrm{I}} 
\newcommand{\rmK}{\mathrm{K}}
\newcommand{\rmM}{\mathrm{M}}
\newcommand{\rmP}{\mathrm{P}}  
\newcommand{\rmp}{\mathrm{p}}  
\newcommand{\rmQ}{\mathrm{Q}}  
\newcommand{\rmR}{\mathrm{R}}
\newcommand{\rmS}{\mathrm{S}}
\newcommand{\rmT}{\mathrm{T}}
\newcommand{\rmU}{\mathrm{U}}
\newcommand{\rmV}{\mathrm{V}}
\newcommand{\rmY}{\mathrm{Y}}
\newcommand{\rmZ}{\mathrm{Z}}
\newcommand{\rmz}{\mathrm{z}}

% arrows and symbols 
\renewcommand{\to}{\rightarrow}
\newcommand{\toto}{\longrightarrow}
\newcommand{\mapstoto}{\longmapsto}
\newcommand{\then}{\Rightarrow}
\newcommand{\IFF}{\Leftrightarrow}
\newcommand{\tl}{\tilde}
\newcommand{\wtl}{\widetilde}
\newcommand{\ol}{\overline}
\newcommand{\ul}{\underline}
\newcommand{\oldemptyset}{\emptyset}
\renewcommand{\emptyset}{\varnothing}
\DeclareMathSymbol{\Arg}{\mathbin}{AMSa}{"39} % for arguments 
\newcommand{\onto}{\ensuremath{\twoheadrightarrow}}
\newcommand{\tle}{\trianglelefteq}
\newcommand{\tge}{\trianglerighteq}

% custom ops
\newcommand{\rmKer}{\mathrm{Ker}}
\newcommand{\rmIm}{\mathrm{Im}}
\newcommand{\lcm}{\mathrm{lcm}}
\newcommand{\GL}{\mathrm{GL}}
\newcommand{\ev}{\mathrm{ev}}
\newcommand{\End}{\mathrm{End}}
\newcommand{\rmid}{\mathrm{id}}
\newcommand{\rmspan}{\mathrm{span}}
\newcommand{\Spec}{\mathrm{Spec}}

% absolute value symbol
\usepackage{mathtools} 
\DeclarePairedDelimiter\abs{\lvert}{\rvert}%
\DeclarePairedDelimiter\norm{\lVert}{\rVert}%
\makeatletter
\let\oldabs\abs
\def\abs{\@ifstar{\oldabs}{\oldabs*}}

% tensor symbol
\newcommand{\tensor}[1]{%
  \mathbin{\mathop{\otimes}\limits_{#1}}%
}

% permutation cycle notation
\ExplSyntaxOn
\NewDocumentCommand{\cycle}{ O{\;} m }
 {
  (
  \alec_cycle:nn { #1 } { #2 }
  )
 }

\seq_new:N \l_alec_cycle_seq
\cs_new_protected:Npn \alec_cycle:nn #1 #2
 {
  \seq_set_split:Nnn \l_alec_cycle_seq { , } { #2 }
  \seq_use:Nn \l_alec_cycle_seq { #1 }
 }
\ExplSyntaxOff

% setminus symbol
\newcommand{\mysetminusD}{\hbox{\tikz{\draw[line width=0.6pt,line cap=round] (3pt,0) -- (0,6pt);}}}
\newcommand{\mysetminusT}{\mysetminusD}
\newcommand{\mysetminusS}{\hbox{\tikz{\draw[line width=0.45pt,line cap=round] (2pt,0) -- (0,4pt);}}}
\newcommand{\mysetminusSS}{\hbox{\tikz{\draw[line width=0.4pt,line cap=round] (1.5pt,0) -- (0,3pt);}}}
\newcommand{\sm}{\mathbin{\mathchoice{\mysetminusD}{\mysetminusT}{\mysetminusS}{\mysetminusSS}}}

% miscellaneous commands
\newcommand{\defn}[1]{{\color{mylightblue}{#1}}}
\newcommand{\toDo}{{\bf\color{red} TODO}}
\newcommand{\toCite}{{\bf\color{green} CITE}}
\newcommand*{\vertbar}{\rule[-1ex]{0.5pt}{2.5ex}} % for matrices with column vectors
\newcommand*{\horzbar}{\rule[.5ex]{2.5ex}{0.5pt}} % for matrices with row vectors
\newcommand{\myblue}[1]{{\color{iceberg}{#1}}}
\newcommand{\myorange}[1]{{\color{burntorange}{#1}}}
\newcommand{\mygreen}[1]{{\color{applegreen}{#1}}}
\newcommand{\myred}[1]{{\color{darkcandyapplered}{#1}}}

% ferrer's diagram
\newcommand{\fdiagram}[1]{
    \begin{tikzpicture}[scale=.7]
        \fill foreach \Z [count=\Y] in {#1}
        {foreach \X in {1,...,\Z} 
        {(\X,-\Y) circle[radius=3pt]}};
    \end{tikzpicture}
}

%
\usepackage{pgfplots}
\pgfplotsset{compat=1.18}

%
\makeatletter
\def\mathcolor#1#{\@mathcolor{#1}}
\def\@mathcolor#1#2#3{%
  \protect\leavevmode
  \begingroup
    \color#1{#2}#3%
  \endgroup
}
\makeatother

% 
\newenvironment{nouppercase}{%
  \let\uppercase\relax%
  \renewcommand{\uppercasenonmath}[1]{}}{}

% titlepage
\title{226: Θεωρία Δακτυλίων και Modules \\ Ασκήσεις με Ενδεικτικές Λύσεις}
\author[Β.~Δ. Μουστακας]{Βασίλης Διονύσης Μουστάκας \\ Πανεπιστήμιο Κρήτης}
\date{4 Μαρτίου 2026}
% \urladdr{\href{https://sites.google.com/view/vasmous}{https://sites.google.com/view/vasmous}}

\begin{document}

\begingroup
\def\uppercasenonmath#1{} % this disables uppercase title
\let\MakeUppercase\relax % this disables uppercase authors
\maketitle
\endgroup

\setcounter{exercise}{4}
% \setcounter{theorem}{5}
\setcounter{equation}{2}
\renewcommand{\thesubsection}{\thesection.\arabic{subsection}}

Σε ότι ακολουθεί υποθέτουμε ότι $R$ είναι ένας \emph{μεταθετικός} δακτύλιος και $n$ είναι ένας θετικός ακέραιος, εκτός αν αναφέρεται διαφορετικά.

\begin{exercise}{\rm(\cite[Άσκηση~7.2.8]{DF04})}
  \label{ex:DF_7.2.8}
  Για κάθε $n \ge 2$, να δείξετε ότι αν $A \in \rmM_n(R)$ είναι αυστηρά άνω τριγωνικός πίνακας (δηλαδή, άνω τριγωνικός πίνακας με $0$ στην κεντρική διαγώνιο), τότε $A^n=0$. Με άλλα λόγια, κάθε αυστηρά άνω τριγωνικός πίνακας είναι μηδενοδύναμο στοιχείο του $\rmM_n(R)$.
\end{exercise}

\begin{proof}[Λύση]
  Αρκεί να αποδείξουμε ότι υψώνοντας τον $A$ σε κάποια δύναμη, ουσιαστικά \textquote{σπρώχνουμε} τα μη μηδενικά στοιχεία του κάποιες διαγωνίους προς τα πάνω. Κατά συνέπεια, μόλις εξαντληθούν οι διαγώνιοι του $A$ ($n$ το πλήθος) προκύπτει ο μηδενικός πίνακας.

  Γι αυτό, θα δείξουμε ότι για κάθε $1 \le k \le n$, αν $A^k = (a_{ij}^{(k)})_{1\le i, j\le n}$, τότε 
  \begin{equation}
    \label{eq:nilpotent_matrix}
    a_{ij}^{(k)} = 0, \quad \text{αν $j < i+k$}.
  \end{equation}
  Θα κάνουμε επαγωγή ως προς $k$. Για $k=1$ η συνθήκη \eqref{eq:nilpotent_matrix} είναι ισοδύναμη με το να είναι ο πίνακας $A$ αυστηρά άνω τριγωνικός (γιατί;), το οποίο ισχύει εξ' υποθέσεως. Αυτή είναι η βάση της επαγωγής. Έστω $k > 1$ και υποθέτουμε ότι η συνθήκη \eqref{eq:nilpotent_matrix} ισχύει για $k$.

  Θεωρούμε τον πίνακα
  \[
  A^{k+1} = A^{k}A = 
  \begin{pmatrix}
                 &              &           & \\
    a_{i1}^{(k)} & a_{i2}^{(k)} & \cdots    & a_{in}^{(k)} \\
                 &              &           & \\
                 &              &           & \\
  \end{pmatrix}
  \begin{pmatrix}
     & & & a_{1j} & \\
     & & & a_{2j} & \\
     & & & \vdots & \\
     & & & a_{nj} & \\
  \end{pmatrix}.
  \]
  To $(i,j)$-στοιχείο του $A^{k+1}$ ισούται με 
  \[
  a_{ij}^{k+1} = a_{i1}^{(k)}a_{1j} + a_{i2}^{(k)}a_{2j} + \cdots + a_{in}^{(k)}a_{nj}.
  \]
  Από την επαγωγική υπόθεση γνωρίζουμε ότι
  \[
  a_{i\ell}^{(k)} = 0, \quad \text{αν $\ell < i+k$}
  \]
  ενώ από την υπόθεση ότι ο $A$ είναι αυστηρά άνω τριγωνικός έπεται ότι 
  \[
  a_{\ell{j}} = 0, \quad \text{αν $j < \ell+1$}. 
  \]
  Συνδυάζοντας αυτές τις δύο συνθήκες παίρνουμε
  \[
  a_{ij}^{k+1}=0, \quad \text{αν $j < \ell+1 < i + (k+1)$},
  \]
  και το ζητούμενο έπεται.
\end{proof}

\begin{exercise}{\rm(\cite[Ασκήσεις~7.1.14~και~7.4.27]{DF04})}
\label{ex:DF_7.1.14_7.4.27}
Έστω $x \in R$ ένα μηδενοδύναμο στοιχείο, δηλαδή $x^m=0$ για κάποιο $m \ge 1$. Να δείξετε τα εξής.
\begin{itemize}
  \item[(1)] $x=0$ ή το $x$ είναι μηδενοδιαιρέτης.
  \item[(2)] Έχουμε $1-x \in R^\times$.
  \item[(3)] Αν $u \in R^\times$, τότε $x+u \in R^\times$.
  \item[(4)] Για κάθε $r \in R$, έχουμε $1 -rx \in R^\times$.
\end{itemize}
\end{exercise}  

\begin{proof}[Λύση]
  Το (4) λύνεται όμοια με το (2). Για τα υπόλοιπα:
  \begin{itemize}
    \item[(1)] Υποθέτουμε ότι $x \neq 0$ και θεωρούμε τον μικρότερο $m > 1$ για τον οποίο $x^m = 0$. Τότε 
    \[
    0 = x^m = \underbrace{x^{m-1}}_{\neq \, 0} \underbrace{x}_{\neq \, 0}
    \]
    και γι αυτό το $x$ είναι μηδενοδιαιρέτης. Διαφορετικά, $x=0$ και το ζητούμενο έπεται.
    \item[(2)] Παρατηρούμε ότι 
    \[
    1 = 1-x^m = (1-x)(1+x+\cdots + x^m)
    \]
    και κατά συνέπεια $1-x \in R^\times$.
    \item[(3)] Έστω $u \in R^\times$. Παρατηρούμε ότι 
    \[
    u + x = u(1 + u^{-1}x) = u(1 - (-u^{-1}x).
    \]
    Από το Ερώτημα (2) και το γεγονός ότι το $R^\times$ είναι κλειστό ως προς τον πολλαπλασιασμό, αρκεί να δείξουμε ότι το $-u^{-1}x$ είναι μηδενοδύναμο. Πράγματι,
    \[
    (-u^{-1}x)^m = (-u)^m\underbrace{x^m}_{=\,0} = 0
    \]
    και το ζητούμενο έπεται.
  \end{itemize}
\end{proof}

\begin{remark}
  Συνδυάζοντας τις Ασκήσεις \ref{ex:DF_7.2.8} και \ref{ex:DF_7.1.14_7.4.27} προκύπτει ότι οι πίνακες της μορφής 
  \[
  \begin{pmatrix}
    \lambda & 1       & 0      & \cdots & 0       & 0 \\
    0       & \lambda & 1      & \cdots & 0       & 0 \\
    \vdots  & \vdots  & \vdots &        & \vdots  & \vdots\\
    0       & 0       & 0      & \cdots & \lambda & 1 \\
    0       & 0       & 0      & \cdots & 0       & \lambda
  \end{pmatrix}
  \]
  είναι αντιστρέψιμοι. Αυτοί είναι γνωστοί ως \emph{στοιχειώδεις πίνακες Jordan} και αποτελούν τα βασικά συστατικά της \emph{κανονικής μορφής Jordan} ενός τετραγωγικού πίνακα.
\end{remark}

\begin{exercise}{\rm(\cite[Άσκηση~7.3.29]{DF04})}
  \label{ex:DF_7.3.29}
  Το σύνολο $\nN(R)$ όλων των μηδενοδύναμων στοιχείων του $R$ ονομάζεται \defn{μηδενοριζικό} (nilradical) του $R$. Να δείξετε ότι το $\nN(R)$ είναι ιδεώδες του $R$.
\end{exercise}

\begin{proof}[Λύση]
  Θα επαληθεύσουμε τον Ορισμό 1.5. 
  \begin{itemize}
    \item Το $\nN(R)$ περιέχει το μηδέν.
    \item Το $\nN(R)$ είναι κλειστό ως προς την πρόσθεση. Πράγματι, έστω $x, y \in \nN(R)$. Τότε $x^m = y^n = 0$, για κάποια $n \ge 1$. Από την Άσκηση 3 (διωνυμικό θεώρημα) έπεται ότι 
    \begin{align*}
    (x+y)^{m+n-1} 
    &= \sum_{k=0}^{n+m-1}\binom{m+n-1}{k} x^ky^{m+n-1-k} \\
    &= \sum_{k=0}^{m-1}\binom{m+n-1}{k}x^k\underbrace{y^{m+n-1-k}}_{= \, 0} \\
    &\quad + \sum_{k=m}^{m+n-1}\binom{m+n-1}{k}\underbrace{x^k}_{=\,0}y^{m+n-1-k} \\
    &= 0.
    \end{align*}
    Άρα, $x+y \in \nN(R)$.
    \item Το $\nN(R)$ είναι κλειστό ως προς πολλαπλασιασμό με στοιχείο του $R$. Πράγματι, έστω $r\in R$ και $x \in \nN(R)$. Τότε, $x^m=0$, για κάποιο $m \ge 1$. Οπότε, 
    \[
    (rx)^m = r^m x^m = 0
    \]
    και γι αυτό $rx \in \nN(R)$.
  \end{itemize}
\end{proof}

\begin{remark}
  Το αποτέλεσμα της Άσκησης \ref{ex:DF_7.3.29} παύει να ισχύει αν ο $R$ δεν είναι μεταθετικός (βλ. \cite[Άσκηση~7.3.31]{DF04}). Για παράδειγμα, στον $R = \rmM_2(\ZZ)$ ισχύει ότι\footnote{Θυμηθείτε τον συμβολισμό της Άσκησης 1.} $E_{21}, E_{12} \in \nN(R)$, ενώ
  \[
  E_{21} + E_{12} = \begin{pmatrix}
    0 & 1 \\
    1 & 0
  \end{pmatrix} \notin \nN(R),
  \]
  διότι 
  \[
  \begin{pmatrix}
    0 & 1 \\
    1 & 0
  \end{pmatrix}^2 = I_2.
  \]
  Με άλλα λόγια, το μηδενοριζικό του $R$ δεν είναι κλειστό ως προς την πρόσθεση.
\end{remark}

\begin{exercise}{\rm(\cite[Άσκηση~7.3.33]{DF04})}
  \label{ex:DF_7.3.33}
  Έστω πολυώνυμο $f(x) = a_0 + a_1x + \cdots + a_nx^n \in R[x]$.
  Να δείξετε ότι το $f(x)$ είναι αντιστρέψιμο στο $R[x]$ αν και μόνο αν $a_0 \in R^\times$ και τα $a_1, a_2, \dots, a_n \in \nN(R)$.
\end{exercise}

\begin{proof}[Λύση]
  Υποθέτουμε ότι $a_0 \in R^\times$ και $a_i \in \nN(R)$ για κάθε $1 \le i \le n$. Επειδή το μηδενοριζικό είναι ιδεώδες (βλ. Άσκηση \ref{ex:DF_7.3.29}), έπεται ότι 
  \[
  a_1x + a_2x^2 + \cdots + a_nx^n \ \in \ \nN(R).
  \]
  Επειδή $a_0 \in R^\times$, από το Ερώτημα (3) της Άσκησης \ref{ex:DF_7.1.14_7.4.27} έπεται ότι 
  \[
  f(x) = a_0 + (a_1x + a_2x^2 + \cdots + a_nx^n) \ \in \ (R[x])^\times
  \]

  Αντιστρόφως, υποθέτουμε ότι $f(x) \in (R[x])^\times$. Τότε 
  \begin{equation}
    \label{eq:nilpotent_polynomial}
    (a_0 + a_1x + \cdots + a_nx^n)(\underbrace{b_0 + b_1x + \cdots + b_mx^m}_{= \, g(x)}) = 1,
  \end{equation}
  για κάποιο $g(x) \in R[x]$. Από την Ταυτότητα \eqref{eq:nilpotent_polynomial} έπεται ότι $a_0b_0 = 1$ και γι αυτό $a_0, b_0 \in R^\times$. Συγκρίνοντας τους όρους $n+m$ και $n+m-1$ στο ανάπτυγμα της Ταυτότητας \eqref{eq:nilpotent_polynomial} προκύπτει ότι 
  \[
  0 = a_nb_m = a_nb_{m-1} + a_{n-1}b_m.
  \]
  Πολλαπλασιάζοντας με $a_n$ και αναδιατάσσοντας τους όρους παρατηρούμε ότι 
  \[
  a_n^2b_{m-1} = \underbrace{a_n^2b_m}_{=\, a_n(\underbrace{a_nb_m}_{=\,0})} - a_{n-1}\underbrace{a_nb_m}_{=\,0} = 0.
  \]
  Επαγωγικά, προκύπτει ότι 
  \[
  a^{k+1}_nb_{m-k} = 0,
  \]
  για κάθε $k$. Ειδικότερα, για $k=m$ έχουμε $a_n^{m+1}b_0 = 0$ και κατά συνέπεια $a_n^{m+1} = 0$ διότι το $b_0$ είναι αντιστρέψιμο. Άρα, $a_n \in \nN(R)$. Επαναλαμβάνοντας το ίδιο επιχείρημα για το πολυώνυμο 
  \[
  f(x) + (-a_nx^n), 
  \]
  το οποίο είναι αντιστρέψιμο από το Ερώτημα (3) της Άσκησης \ref{ex:DF_7.1.14_7.4.27}, προκύπτει ότι $a_{n-1} \in \nN(R)$. Συνεχίζουμε έτσι μέχρι να φτάσουμε στον συντελεστή $a_1 \in \nN(R)$ και το ζητούμενο έπεται.
\end{proof}

\begin{exercise}{\rm(βλ. \cite[Άσκηση~7.3.33]{DF04})}
  \label{ex:nilradicals}
  Να περιγράψετε το μηδενοριζικό καθενός από τους παρακάτω δακτυλίους:
  \begin{itemize}
    \item[(1)] μιας ακέραιας περιοχής (και ειδικότερα του $\ZZ$),
    \item[(2)] του $\ZZ_n$, και 
    \item[(3)] του πολυωνυμικού δακτυλίου μίας πεταβλητής $R[x]$.
  \end{itemize}
\end{exercise}

\begin{proof}[Λύση]
  \leavevmode
  \begin{itemize}
    \item[(1)] Από το Ερώτημα (1) της Άσκησης \ref{ex:DF_7.1.14_7.4.27} έπεται ότι κάθε μη μηδενικό μηδενοδύναμο στοι\-χείο είναι μηδενοδιαιρέτης. Κατά συνέπεια, το μηδενοριζικό μιας ακέραιας περιοχής είναι το τετριμμένο ιδεώδες. Ειδικότερα, 
    \[
    \nN(\ZZ) = \{0\}.
    \]

    \item[(2)] Στην Άσκηση 4 είδαμε ότι $x \in \nN(\ZZ_n)$ αν και μόνο αν κάθε πρώτος διαιρέτης του $n$ διαιρεί και το $x$. Συνεπώς, αν $p_1, p_2, \dots, p_\ell$ είναι οι πρώτοι που εμφανίζονται στην παραγοντοποίηση του $n$ σε πρώτους, τότε 
    \[
    \nN(\ZZ_n) = p_1\ZZ_n \cap p_2\ZZ_n \cap \cdots \cap p_\ell\ZZ_n = (p_1p_2\cdots{p_\ell})\ZZ_n.
    \]

    \item[(3)] Έστω $f(x) = a_0 + a_1x + \cdots + a_nx^n \in R[x]$. Θα δείξουμε ότι 
    \[
    \nN\left(
      R[x]
    \right) = \nN(R)[x].
    \]
    Υποθέτουμε ότι $a_i \in \nN(R)$, για κάθε $1 \le i \le n$. Αυτό σημαίνει ότι $a_ix^i \in \nN\left(R[x]\right)$, για κάθε $1 \le i \le n$ (γιατί;) και κατά συνέπεια $f(x) \in \nN\left(R[x]\right)$ (γιατί;). Άρα, $\nN(R)[x] \subseteq \nN\left(R[x]\right)$.

    Αντιστρόφως, αν $f(x) \in \nN\left(R[x]\right)$, τότε $xf(x) \in \nN\left(R[x]\right)$ (γιατί;) και γι αυτό 
    \[
    1 + xf(x) \ \in \ \left(R[x]\right)^\times
    \]
    από το Ερώτημα (3) της Άσκησης \ref{ex:DF_7.1.14_7.4.27}. Από την Άσκηση \ref{ex:DF_7.3.33}, έπεται ότι $f(x) \in \nN(R)[x]$ (γιατί;). Συνεπώς, $\nN\left(R[x]\right) \subseteq \nN(R)[x]$ και η λύση ολοκληρώνεται.
  \end{itemize}
\end{proof}

\begin{remark}
  Με τον συμβολισμό που αναπτύξαμε, ο υπολογισμός της Άσκησης 4 μπορεί να γραφεί ως 
  \[
  \nN(\ZZ_{72}) = 2\ZZ_{72} \cap 3\ZZ_{72} = 6\ZZ_{72}.
  \]
\end{remark}

\begin{exercise}{\rm(\cite[Άσκηση~7.4.26]{DF04})}
  Έστω $\Spec(R)$ το σύνολο όλων των πρώτων ιδεωδών του $R$. Να δείξετε ότι 
  \[
  \nN(R) \subseteq \bigcap_{\fp \in \Spec(R)} \fp.
  \]
\end{exercise}

\begin{proof}[Λύση]
  Αρκεί να δείξουμε ότι κάθε πρώτο ιδεώδες του $R$ περιέχει κάθε μηδενοδύναμο στοιχείο του $R$ (γιατί;). Πράγματι, έστω $\fp \in \Spec(R)$ και $x \in \nN(R)$. Τότε $x^m = 0$, για κάποιο $m \ge 1$. Συνεπώς,
  \[
  0 = x^m \ \in \ \fp 
  \]
  και γι αυτό $x \in \fp$, διότι το $\fp$ είναι πρώτο. Άρα, $\nN(R) \subseteq \fp$ και το ζητούμενο έπεται.
\end{proof}

\begin{thebibliography}{99}
    \bibitem{DF04}
    D. S. Dummit και R. M. Foote, 
    \textit{Abstract algebra},
    third edition, Wiley, 2004. 
\end{thebibliography}
\end{document}